% \vspace{0.20 cm}
\section{Research Experience}
        \vspace{0.10 cm}
        \begin{twocolentry}{
            Nov 25 - Present
        }
            \textbf{Study of X-Ray Flares and Reverse Shock Mechanism in GRB Prompt Phase.} \end{twocolentry}

        %\vspace{0.10 cm}
        \begin{onecolentry}
            \textbf{\textit{Collaborators: Dr. Bin Bin Zhang (NJU), Dr. Johnathen Granot (OUI and GWU), Zhen-Yu Yan (NJU).}}\\
            
            \vspace{0.10 cm}
            \begin{highlights}
                \item Investigating historical X-ray and Gamma-ray archives to identify early-time \textbf{\textit{X-ray flares}}, motivated by recent observations from the \textbf{\textit{Einstein Probe mission}}.
            \item Implementing the two-parameter \textbf{\textit{reverse shock model}} proposed by Granot et al. to characterize the physics of the GRB prompt emission phase.
            \end{highlights}

        \end{onecolentry}
        
        \vspace{0.10 cm}
        \begin{twocolentry}{
            July 25 - Present
        }
            \textbf{Multiwavelength analysis of KM3-230 213A Event.} \end{twocolentry}

        %\vspace{0.10 cm}
        \begin{onecolentry}
            \textbf{\textit{Colaborators: Dr. Jagdish C. Joshi (ARIES), Dr. Timur Dzhatdoev (INR RAS).}}\\
            
            \vspace{0.10 cm}
            \begin{highlights}
                \item Carrying out a multi-messenger analysis of the most energetic neutrino event, \textbf{\textit{KM3-230213A}}, detected by KM3NeT, by incorporating data from \textbf{\textit{UHECRs}}, \textbf{\textit{neutrinos}}, and \textbf{\textit{$\gamma$-ray}} observations and using \textbf{\textit{CRPropa}} simulations.
            \end{highlights}

        \end{onecolentry}
        
        \vspace{0.10 cm}
        \begin{twocolentry}{
            Dec 24 - Nov 24
        }
            \textbf{Study of VHE Component in Gamma Ray Bursts and Pulsar Wind Nebulae.} \end{twocolentry}

        %\vspace{0.10 cm}
        \begin{onecolentry}
            \textbf{\textit{Guide : Dr. Jagdish C. Joshi, Aryabhatta Research Institute of Observational Sciences (ARIES), India.}}\\
            
            \vspace{0.10 cm}
            \begin{highlights}
                \item Conducted studies on the \textbf{\textit{minimum variability timescales (MVT)}} of GRBs, explored parameter degeneracies in models using the \textbf{\textit{MCMC}} algorithm, and investigated non-thermal VHE emission through radiative processes such as Synchrotron, IC (Inverse Compton), and SSC (Synchrotron Self-Compton) during both the prompt and afterglow phases of GRBs. We also performed analytical calculations within the \textbf{\textit{Synchrotron-SSC}} framework to examine the \textbf{\textit{VHE}} (GeV-scale) component.
                
                \item We performed some analytical calculations on an alternative mechanism for the VHE component, namely \textbf{\textit{Jitter Radiation}}, in the context of GRBs and PWNe. In addition, we are investigating the two-hump spectrum observed in \textbf{\textit{GRB 240825A}} using \textbf{\textit{LeHaMoC}} (Lepto-Hadronic-Modeling-Code).


                % and \textbf{\textit{$\gamma \pi$}} (GammaPy) for this study.
                % \item \textbf{Skills Developed :} Python, GRB data analysis, MCMC algorithms (bilby, PyMultinest, emcee), TensorFlow.
                % \item Studied \textbf{\textit{minimum variability time scales (MVT)}} in short and long GRBs.
                % \item Verified results for non-degenerate and degenerate GRB afterglow model like (\textbf{\textit{GRB170817A}}) using \textbf{\textit{Markov Chain Monte Carlo (MCMC)}} sampling algorithms.
                % \item Focusing on the causes and effects of parameter degeneracy in afterglow model and the challenges it presents while fitting the data to the model, particularly when using the MCMC method. Exploring alternate methods to MCMC for interpreting the data to reduce and overcome parameter degeneracy.
                % \item Studied the \textit{\textbf{synchrotron mechanism}} and \textit{\textbf{SSC component}} typical very high energy (VHE) GRBs like \textit{\textbf{GRB 080916C}}.
                
            \end{highlights}        
        \end{onecolentry}
        
        \vspace{0.2 cm}
        \begin{twocolentry}{
            Aug 24 - June 25
        }
            \textbf{Study of Presence of Dark Matter - Core in Different Compact Stellar Models} \end{twocolentry}

        \vspace{0.10 cm}
        \begin{onecolentry}
            \textbf{\textit{Guide : Dr. Apratim Ganguly, Inter-University Centre of Astronomy and Astrophysics (IUCAA), India.}}\\
            
            \vspace{0.10 cm}
            \begin{highlights}
                \item Studied compact stars using admixed models like \textbf{\textit{two-fluid}} and \textbf{\textit{FBS}} (Fermion-Boson Star) with dark matter core configuration by solving \textbf{\textit{TOV}} (Tolman-Oppenheimer-Volkoff), \textbf{\textit{EKG}} (Einstein-Klein-Gordon) and \textbf{\textit{EHKG}} (Einstein-Hilbert-Klein-Gordon) equations numerically.
                
                \vspace{0.10 cm}
                
                \item Studied effect of fraction of dark matter on physical parameters like M (total gravitational mass), M-R (mass-radius relation), $\lambda$ (tidal deformability), C (compactness) etc. \textit{(publication in preparation)}

                % \item \textbf{Skills Developed :} Python, GW data analysis, Julia, Mathemtica, Git. 
                % \item Modeled \textit{\textbf{compact objects}}, including neutron stars, boson stars, and two-fluid admixed models of compact star with dark matter and ordinary matter.

                % \item Solved numerically, the \textit{\textbf{Tolman-Oppenheimer-Volkoff (TOV)}} equations for neutron stars and \textit{\textbf{Einstein-Klein-Gordon (EKG)}} equations for boson stars using advanced numerical methods and developed algorithms for the same.

                % \item Analyzing \textit{\textbf{Fermion Boson Star (FBS)}} and \textit{\textbf{two-fluid admixed model}} to constrain dark matter fractions in compact stars using gravitational wave \textbf{\textit{LIGO}} data, such as from \textit{\textbf{GW170817}} observation.

                % \item Studied compact star properties (M-R relation, tidal deformability, compactness, dark matter fraction etc.), with results documented in a research paper (\textit{in preparation}).
                
            \end{highlights}        
        \end{onecolentry}
        
        \vspace{0.2 cm}
        \begin{twocolentry}
            {Jan - Aug 24}
            
            \textbf{Optimization Of The Feature Tracking Code To Analyze Solar Quiet Sun Data \href{https://drive.google.com/file/d/177hjcXcHCzlMHXnNjdPv3_OTtgu7IE0A/view?usp=sharing}{\faExternalLink}}{\href{https://drive.google.com/file/d/1XAY01jkwndByMHBNCZhTZ_OLq4SwbaKe/view?usp=sharing}{(Masters Thesis)}}    
            
        \end{twocolentry}
        \vspace{0.10 cm}
        \begin{onecolentry}
            \textbf{\textit{Guide : Dr. Anusha L. S., Indian Institute of Astrophysics (IIA), India.}}\\

            \vspace{0.10 cm}
            \begin{highlights}
                \item Converted \textit{\textbf{IDL}} code by \href{https://arxiv.org/abs/1608.08499}{Anusha et. al. (2017)}  into \textit{\textbf{Python}} and \textit{\textbf{Julia}} for open-source availability and cross-platform compatibility. Developed efficient pipelines for the feature trcaking code and optimized it. 
                
                \vspace{0.10 cm}
                
                \item Tested with \textit{\textbf{SUNRISE-II}} data by \textit{\textbf{MPS}} showing that runtime of code is reduced from \textit{\textbf{7 hours}} to \textit{\textbf{2.5 hours}} in Python and from \textit{\textbf{1 hour}} to \textit{\textbf{25 minutes}} in Julia, surpassing the IDL runtime of \textit{\textbf{40 minutes}}. Code is being used for analyzing data from \textit{\textbf{SUNRISE III}} observations. \textit{(publication in preparation)}

                
                % \item \textbf{Skills Developed :}  Python, Julia, Git, Linux machines (servers), Parallelization. 
                % \item Analyzed small-scale magnetic features on the solar photosphere in the quiet Sun region using \textit{\textbf{Sunrise Observatory data by Max Planck Institute for Solar System Research (MPS)}}.
                % \item Converted \textit{\textbf{IDL}} code by \href{https://arxiv.org/abs/1608.08499}{Anusha et. al. (2017)}  into \textit{\textbf{Python}} and \textit{\textbf{Julia}} for open-source availability and cross-platform compatibility.
                % \item Developed \textit{\textbf{efficient pipelines}} for the code and \textit{\textbf{optimized}} the code using (loop optimization, efficient data structures, parallelization, and just-in-time compilation), reducing runtime from \textit{\textbf{7 hours}} to \textit{\textbf{2.5 hours}} in Python and from \textit{\textbf{1 hour}} to \textit{\textbf{25 minutes}} in Julia, surpassing the IDL runtime of \textit{\textbf{40 minutes.}}
                % \item Code is being used for analyzing data from \textit{\textbf{SUNRISE III}} observations with results documented in a research paper (\textit{in preparation}).
                % \item Runtime of the code without parallelization is given in following table
                
                % \begin{center}
                %    \begin{tabular}{|c|c|c|}
                %     \hline
                %     \textbf{Programming Language} & \textbf{Unoptimized Runtime}& \textbf{Optimized Runtime}\\
                %     \hline
                %     IDL & 1 hr & 40 min \\ 
                %     \hline
                %     Python & 7 hr & 2 hr 30 min \\ 
                %     \hline
                %     Julia & 1 hr & 25 min \\ 
                %     \hline
                %    \end{tabular}
                % \end{center}
            \end{highlights}
            
        \end{onecolentry}
        
